\documentclass[12pt,a4paper,italian]{article}


\usepackage[italian]{babel}
\usepackage[latin1]{inputenc}
\usepackage{amsmath}
\usepackage{amsfonts}
\usepackage{amssymb}
\usepackage{color}
\usepackage{xcolor}
\usepackage{hyperref}
\usepackage[all]{hypcap}
\usepackage{ifthen}
\usepackage{wrapfig}
%\author{Piero Bizzotto}
\usepackage[top=2cm,bottom=5cm,left=80pt,right=80pt]{geometry}
\usepackage{graphicx}
\DeclareGraphicsExtensions{.jpg,.png}

\newcommand{\ajax}{AJAXDRAW}
\newcommand{\sito}{\href{http://ajaxdraw.sourceforge.net}{http://ajaxdraw.sourceforge.net}}

\setlength{\parindent}{0pt} %settato indentazione di default 
\setlength{\headheight}{3cm} %settato grandezza header...in altre parole, quanto distanzio il doc dall'intestazione

\usepackage{fancyhdr} %pacchetto per le intestazioni
\pagestyle{fancy} %uso del pacchetto


\fancyhead{} %annulla head di default
\fancyfoot{} %annulla foot di default


\usepackage{lastpage} %setto pg di pgtot a rfoot
     \rfoot{pagina \thepage\ di \pageref{LastPage}}


\lfoot{Versione: \insertversion} %setto versione doc a lfoot
\renewcommand{\footrulewidth}{0.5pt} %ridefinisco il valore della riga di intestazione
\renewcommand{\headrulewidth}{0.5pt} %ridefinisco il valore della riga di pie' di pagina

\newcommand{\insertversion}{0.0} %definisco il nuovo comando per inserire la versione


\lhead{  \begin{Huge} \ajax \end{Huge} \\  %intestazione di sinistra
					%\begin{Large}	Software per il Disegno Grafico\\ in Tecnologie Web \end{Large}  
					\begin{normalsize}\sito \end{normalsize}
			%\\ versione documento: \insertversion\ del \today} %setto l'intestazione sx
		}
\rhead{ %includo logo nell'intestazione dx
	 	\includegraphics[scale=0.5]{../logo/logo.png}  
}


%CREAZIONE ELENCHI NUMERATI PERSONALIZZATI
\newcounter{Lcount}
\newcounter{Rcount}
\setcounter{Lcount}{0}
\setcounter{Rcount}{0}

\newenvironment{elenconumerato}[2][ ]
{
  \begin{list}{#1\arabic{Lcount}.}
    {
	\setcounter{Rcount}{\value{Lcount}}
	\setcounter{Lcount}{0} 
	\usecounter{Lcount} 
\addtolength{\leftmargin}{#2pt}
	}
}
{
  \end{list}
 \setcounter{Lcount}{\value{Rcount}}
}

%CREAZIONE ELENCHI PUNTATI
\newenvironment{elencopuntato}[1][]
{
\begin{list}{\textbullet} %\itemindent=#1pt
	{
	\addtolength{\leftmargin}{#1pt}
	}
} 
{
\end{list}
}


\newenvironment{elencodescrittivo}[1][]{\begin{description} \setlength{\itemindent}{#1pt} \addtolength{\leftmargin}{#1pt}} {\end{description}}

\newcommand{\TITOLODOC}{Titolo}

%footer centrale
\cfoot{ \TITOLODOC \\  E-mail:    \href{ mailto:webshape.contact@gmail.com}{ webshape.contact@gmail.com}  }

%INSERIMENTO IMMAGINI
\newcommand{\imagerealsize}[1]{\vspace{20pt} \includegraphics{#1} }
\newcommand{\imageadapted}[1]{\vspace{20pt} \includegraphics[width=1\textwidth]{#1} }

\newcommand{\glosspath}{.\glossario}
\newcommand{\gloss}[1]{\hyperref{\glosspath~\glossario.pdf}{}{#1}{#1}}

\hypersetup{
    %bookmarks=true,         % show bookmarks bar?
    %unicode=false,          % non-Latin characters in Acrobat’s bookmarks
	%pdftoolbar=true,        % show Acrobat’s toolbar?
	%pdfmenubar=true,        % show Acrobat’s menu?
    %pdffitwindow=true,      % page fit to window when opened
    %pdftitle={My title},    % title
    %pdfauthor={Author},     % author
    %pdfsubject={Subject},   % subject of the document
    %pdfnewwindow=true,      % links in new window
    %pdfkeywords={keywords}, % list of keywords
    colorlinks=true,         % false: boxed links; true: colored links
    linkcolor=black,           % color of internal links
    %citecolor=green,        % color of links to bibliography
    %filecolor=magenta,      % color of file links
    urlcolor=teal    % color of external links
%	linktocpage=false;
}


%COLORAZIONE TESTO
\newcommand{\blue}[1]{{\color {blue} #1}} 
\newcommand{\red}[1]{{\color {red} #1}}
\newcommand{\green}[1]{{\color {green} #1}}
\newcommand{\sezione}[1]{\leftskip=0pt \section{#1} \leftskip=18pt}
\newcommand{\subsezione}[1]{\leftskip=18pt \subsection{#1} \leftskip=36pt}
\newcommand{\subsubsezione}[1]{\leftskip=36pt \subsubsection{#1} \leftskip=54pt}
\newcommand{\subsubsecindent}{54}
\newcommand{\subsecindent}{36}
\newcommand{\secindent}{18}
\newcommand{\normindent}{8}
\newcommand{\code}[1]{{\bfseries \texttt{#1}}}
\newcommand{\paragrafo}[1]{\leftskip=36pt \paragraph{#1} \leftskip=54pt}
\newcommand{\subparagrafo}[1]{\leftskip=54pt \subparagraph{#1} \leftskip=72pt} %BASE!!!
\usepackage{multirow}
\title{\TITOLODOC}
\author{Marco Cunico}

\begin{document}

\renewcommand{\insertversion}{0.0} %INSERIRE LA VERSIONE QUI DENTRO STILE x.x.xx
\renewcommand{\TITOLODOC}{Manuale Utente} %INSERIRE IL TITOLO DEL DOCUMENTO DA FAR COMPARIRE A PIE PAGINA
\renewcommand{\glosspath}{.\glossario} %INSERIRE PERCORSO RELATIVO

%%%%%%%%%%%%%%%%%%%%%%PARTE DA NON MODIFICARE%%%%%%%%%%%%%%%%%
\begin{titlepage}
\begin{center}
	\begin{Large}	\today \end{Large}
\end{center}

\vspace{20pt}

\begin{center}
	\begin{Huge}
				\textbf{\ajax}
	\end{Huge}
\end{center}			

\begin{center}
	\begin{large}
				\textbf{Software per il Disegno Grafico\\ in Tecnologie Web}
	\end{large}
\end{center}			

\vspace{20pt}

\begin{center}
\includegraphics[width=150pt]{../logo/logo}
\end{center}

\vspace{170pt}
\begin{center} %INSERIRE ALL'INTERNO IL TITOLO DOCUMENTO CHE COMPARIRA NELLA PAGINA INIZIALE				
	\begin{Huge}
				\textbf{\TITOLODOC}
	\end{Huge}
			\\
\end{center}
\vspace{190pt}
\begin{center}
Versione: \insertversion
\end{center}
\end{titlepage}

\newpage
%%%%%%%%%%%%%%%%%%%%%%FINE PARTE DA NON MODIFICARE%%%%%%%%%%%%%%%%%

\begin{center} %INSERIRE ALL'INTERNO IL TITOLO DOCUMENTO CHE COMPARIRA NELLA PAGINA INIZIALE
	\begin{Huge}	
				\textbf{\TITOLODOC}
			\\
	\end{Huge}
\end{center}

%\setlength{\parindent}{18pt} %settato indentazione di default 
\section*{\LARGE Sommario:}
Il presente documento descrive l'analisi dei costi e delle risorse effettuata dall'azienda WebShape per lo sviluppo del Capitolato C04.

 %SEZIONE SOMMARIO
\indent \indent

\section*{\LARGE Stato del documento:}
\indent \indent
	Formale Esterno

\section*{\LARGE Redazione:}
	\begin{table}[!h]
		\begin{center}
			\begin{tabular}
				{|c|c|}
				\hline
				%%%%%%%%%%%%%%INTESTAZIONE COLONNE%%%%%%%%%%%%%%%%%%%%%%%%%%%%%%%%
				\multicolumn{2}{|c|}{ \textbf{Redazione} } \\
				\hline
				\textbf{Fase} & \textbf{Redattori} \\
				%%%%%%%%%%%%%%FINE INTESTAZIONE COLONNE%%%%%%%%%%%%%%%%%%%%%%%%%%%%%%%%%%%%%%
				\hline
				%%%%%%%%%%% PARTE DA MODIFICARE %%%%%%%%%%%%%%%%%%%%%%%%%%%%%%%%%%%%%%%%%%		
				\multirow{2}{*}{PPP-RQ} & \\
										& \\
				\hline
				%%%%%%%%%%% FINE PARTE DA MODIFICARE %%%%%%%%%%%%%%%%%%%%%%%%%%%
			\end{tabular}
			\caption{Lista Redattori} %INSERIRE DIDASCALIA - SE NECESSARIA - 
			\label{tabredazione}
		\end{center}
	\end{table}
	
	
\section*{\LARGE Approvazione:}
\begin{table}[!h]
	\begin{center}
		\begin{tabular}
			{|c|c|}
			\hline
			%%%%%INTESTAZIONE COLONNE%%%%%%%%%%%%%%%%%%%%%%%%%%%%%%%
			\multicolumn{2}{|c|}{ \textbf{Approvazione} } \\
			\hline
			\textbf{Fase} & \textbf{Approvatori} \\
			%%%%%%%%%%%%%%FINE INTESTAZIONE COLONNE%%%%%%%%%%%%%%%%%%%%%%%%%%%%%%
			\hline
			%%%%%%%%%%% PARTE DA MODIFICARE %%%%%%%%%%%%%%%%%%%%%%%%%%%%%%%%%%%%%%		
			\multirow{1}{*}{RPP-RQ} & \\
									
			\hline
			%%%%%%%%%%% FINE PARTE DA MODIFICARE %%%%%%%%%%%%%%%%%%%%%%%%%%%%%%%%%%%
		\end{tabular}
		\caption{Lista Approvatori} %INSERIRE DIDASCALIA - SE NECESSARIA - 
		\label{tabapprovazione}
	\end{center}
\end{table}

\textbf{}
\newpage
\section*{\LARGE Lista di Distribuzione:}

	\begin{elenconumerato}{\normindent}
		\item WebShape 
		\item I committenti Conte Renato e Vardanega Tullio in rappresentanza \\  dell'azienda proponente Zucchetti SPA
	\end{elenconumerato}




\section*{\Large Registro delle Modifiche:}


\begin{center}
	\begin{table}[h]
		  \begin{tabular*}
			{1\textwidth}%
				{@{\extracolsep{\fill}}|p{0.1\textwidth}|p{0.54\textwidth}|p{0.26\textwidth}|}
			 \hline
%%%%%%%%%%%%%%INTESTAZIONE COLONNE%%%%%%%%%%%%%%%%%%%%%%%%%%%%%%%%%%%%%%%%%%
			\textbf{Versione}  & \textbf{Descrizione} & \textbf{Autore} \\
%%%%%%%%%%%%%%FINE INTESTAZIONE COLONNE%%%%%%%%%%%%%%%%%%%%%%%%%%%%%%%%%%%%%%%
		 \hline
%%%%%%%%%%% PARTE DA MODIFICARE %%%%%%%%%%%%%%%%%%%%%%%%%%%%%%%%%%%%%%%%%%%
    	  \hline
    	  0.0 & 27/02/2008 Prima stesura & \\

		\hline %%FINE RIGA
%%%%%%%%%%% FINE PARTE DA MODIFICARE %%%%%%%%%%%%%%%%%%%%%%%%%%%%%%%%%%%%%
		\end{tabular*}
	\caption{Registro delle modifiche} %INSERIRE DIDASCALIA - SE NECESSARIA - 
	\label{tab:modifiche}
	\end{table}
\end{center}


\newpage
\thispagestyle{fancy}
\tableofcontents
\thispagestyle{fancy}
\newpage

\sezione{Introduzione}

\subsezione{Scopo del documento}
Il presente documento si propone di fornire la documentazione di supporto necessaria per l'utilizzo del prodotto.\\

\subsezione{Definizione dell'utente del prodotto}
Il prodotto \`e pensato per tutti gli utenti che hanno conoscenze minime in ambiente di grafica vettoriale e che possiedono una connessione ad internet. Il sistema potr\`a anche essere usato da chiunque voglia estendere l'applicazione, secondo le norme che regolano la licenza d'utilizzo.\\

\subsezione{Come leggere il manuale}
Il presente Manuale Utente \`e suddiviso in varie sezioni. Nella sezione ''Introduzione'' il lettore pu\`o trovare riferimenti a documenti utili per l'uso del prodotto e informazioni riguardo eventuali segnalazioni di problemi riscontrati nel suo utilizzo.
Nella sezione ''Descrizione generale'' viene descritto il prodotto illustrandone lo scopo, le modalit\`a di utilizzo, le caratteristiche salienti, i requisiti tecnici. Nella sezione ''Istruzioni d'uso'' sono presenti le operazioni da effettuare per utilizzare il software, a partire dall'installazione fino al termine di una sessione di utilizzo.\\ %integrate anche mediante vari screenshot


\subsezione{Documenti Utili}
\begin{elencopuntato}[\normindent]
	\item[-] \textit{Glossario.pdf}
\end{elencopuntato}

\subsezione{Come riportare problemi e malfunzionamenti}
Nel caso durante l'utilizzo del prodotto sorgano problemi o malfunzionamenti non descritti nel seguente documento, l'utente \`e invitato a contattare la WebShape all'indirizzo \href{mailto:webshape.contact@gmail.com}{webshape.contact@gmail.com} riportando i seguenti dati:\\
\begin{elencopuntato}[\normindent]
	\item[-] Versione del prodotto
	\item[-] Sistema operativo in uso
	\item[-] Browser in uso
	\item[-] Configurazione hardware in uso
	\item[-] Fase di utilizzo in cui si \`e manifestato l'errore
	\item[-] Descrizione dell'errore
\end{elencopuntato}

\sezione{Descrizione Generale}

\subsezione{Caratteristiche del Prodotto}
Il prodotto \`e stato sviluppato come un'applicazione web. \`E possibile utilizzarlo, una volta installato sul server, semplicemente accedendovi tramite un browser. Una volta effettuato l'accesso alla pagina \`e possibile accedere a tutte le
funzionalit\`a. Il sistema permette un salvataggio dei disegni che potranno essere ricaricati e aggiornati.

\subsezione{Requisiti tecnici per il funzionamento del programma}
\begin{elencopuntato}[\normindent]
    \item[-] Requisiti software:  Elenco browser supportati, risoluzione ecc..
    \item[-] Requisiti hardware:  Processore .. memoria RAM ..
\end{elencopuntato}

\sezione{Istruzioni per l'uso}

\subsezione{Installazione del programma}
Il programma verr\`a fornito come?? con una preinstallazione, pacchettizzato, gi\`a pronto all'installazione??.

\subsezione{Avvio del programma}

\subsezione{Descrizione funzionale}

\sezione{Appendice}

\subsezione{Messaggi di errore e loro significato}
\end{document}
